\documentclass{article}

\usepackage[utf8]{inputenc}     % pdfLaTeX
\usepackage[T1]{fontenc}
\usepackage[magyar]{babel}
\usepackage{amsmath}

% Set page size and margins
% Replace `letterpaper' with `a4paper' for UK/EU standard size
\usepackage[a4paper,top=2cm,bottom=2cm,left=3cm,right=3cm,marginparwidth=1.75cm]{geometry}
%for arrows
\usepackage{textcomp}

\title{Díjak, kitüntetések helyesírása}
\author{Lengyel Zoltán Péter \thanks{Csukás Ibolya tanítása alapján}}
\date{Legutóbb frissítve: 2026. március 16.}

\begin{document}

\maketitle

\tableofcontents

\newpage

\section{Tulajdonnévi előtag}
\begin{itemize}
    \item Az előtagot nagy-, az utótagot kisbetűvel írjuk, és kötőjellel kapcsoljuk össze őket
    \begin{itemize}
        \item Kossuth-díj
        \item Nobel-díj
        \item Oscar-díj
        \item Eötvös Loránd-emlékérem
        \item Alföld-díj
    \end{itemize}

    \item Kiegészítéssel: 
    \begin{itemize}
        \item pl. Nobel-díj-átadás
    \end{itemize}
\end{itemize}

\subsection{-i, -s képzős alak}
\begin{itemize}
    \item Az előtagot nagybetűvel, az utótagot kisbetűvel írjuk \begin{itemize}
        \item pl. Nobel-díjas, Kossuth-díjas
    \end{itemize}
\end{itemize}

\section{Köznévi tagokból álló}
\begin{itemize}
    \item Minden tagot nagybetűvel írunk
    \begin{itemize}
        \item Állami Díj
        \item Nyelvi Díj
        \item Innovációs Díj
    \end{itemize}
\end{itemize}
    \subsection{-i, -s képzős alakok}
    \begin{itemize}
        \item Minden tagot kisbetűvel írunk, a képzőt közvetlenül csatoljuk
        \begin{itemize}
            \item pl. állami díjas, európai nyelvi díjas
        \end{itemize}
    \end{itemize}



\end{document}